% !TEX program = xelatex
\documentclass[11pt,a4paper]{ctexart}

% ==================== 颜色 ====================
\usepackage{xcolor}
\definecolor{maincolor}{HTML}{1A7A6B}       % 主色（标题、顶栏、框边框）
\definecolor{accent}{HTML}{D4A843}          % 强调色（关键词高亮）
\definecolor{boxbg}{HTML}{F0F7F5}           % 问题框背景
\definecolor{borderblue}{HTML}{2E86AB}      % 方法框边框
\definecolor{proofbg}{HTML}{FDF8F0}         % 方法框背景
\definecolor{summarybg}{HTML}{F5F0FF}       % 总结框背景

% ==================== 页面 ====================
\usepackage[top=2.2cm, bottom=2.2cm, left=2.5cm, right=2.5cm, headheight=23.12pt]{geometry}

% ==================== 字体 ====================
\setmainfont{TeX Gyre Pagella}

% ==================== 数学 ====================
\usepackage{amsmath,amssymb,amsfonts,mathtools}

% ==================== 彩色框 ====================
\usepackage{tcolorbox}
\tcbuselibrary{skins,breakable,most}

% ==================== 页眉页脚 ====================
\usepackage{fancyhdr}
\pagestyle{fancy}
\fancyhf{}
\fancyhead[L]{\color{maincolor}\sffamily\bfseries 数学讲义}
\fancyhead[R]{\color{gray}\sffamily \today}
\fancyfoot[C]{\color{gray}\sffamily\thepage}
\renewcommand{\headrulewidth}{1.5pt}
\renewcommand{\headrule}{\hbox to\headwidth{\color{maincolor}\leaders\hrule height \headrulewidth\hfill}}

% ==================== 章节标题样式 ====================
\usepackage{titlesec}
\titleformat{\section}
  {\color{maincolor}\sffamily\LARGE\bfseries}
  {\thesection}{1em}{}
  [\vspace{0.3cm}{\color{maincolor!40}\titlerule[2pt]}]

\titleformat{\subsection}
  {\color{maincolor!80}\sffamily\Large\bfseries}
  {\thesubsection}{1em}{}

% ==================== 表格 ====================
\usepackage{array,booktabs,colortbl}

% ==================== 杂项 ====================
\usepackage{enumitem}
\setlist{nosep, leftmargin=*}
\usepackage[hidelinks]{hyperref}

% ==================== 自定义命令 ====================
\newcommand{\highlight}[1]{\colorbox{accent!20}{\sffamily\bfseries #1}}

% ==================== 自定义环境 ====================

% 1. 问题框
\newtcolorbox[auto counter]{problembox}[1][]{
  enhanced,
  colframe=maincolor,
  colback=boxbg,
  coltitle=white,
  fonttitle=\sffamily\bfseries,
  title={\large 问题},
  attach boxed title to top left={xshift=4mm, yshift=-2mm},
  boxed title style={colback=maincolor, sharp corners},
  sharp corners,
  boxrule=1.2pt,
  left=4mm, right=4mm, top=5mm, bottom=4mm,
  before skip=1.2em, after skip=1.2em,
  #1
}

% 2. 方法/证明框
\newtcolorbox[auto counter]{methodbox}[2][]{
  enhanced,
  colframe=borderblue,
  colback=proofbg,
  coltitle=white,
  fonttitle=\sffamily\bfseries,
  title={#2},
  attach boxed title to top left={xshift=4mm, yshift=-2mm},
  boxed title style={colback=borderblue, sharp corners},
  sharp corners,
  boxrule=1.2pt,
  left=4mm, right=4mm, top=5mm, bottom=4mm,
  before skip=1.2em, after skip=1.2em,
  #1
}

% 3. 总结框
\newtcolorbox{summarybox}[1][]{
  enhanced,
  colframe=maincolor!60,
  colback=summarybg,
  coltitle=white,
  fonttitle=\sffamily\bfseries,
  title={\large $\blacktriangleright$ 总结},
  attach boxed title to top left={xshift=4mm, yshift=-2mm},
  boxed title style={colback=maincolor!70, sharp corners},
  sharp corners,
  boxrule=1.2pt,
  left=4mm, right=4mm, top=5mm, bottom=4mm,
  before skip=1.5em, after skip=1.2em,
  #1
}

% ==================== 文档开始 ====================
\begin{document}

% ---- 彩色顶栏标题 ----
\begin{tcolorbox}[
  enhanced,
  colframe=maincolor,
  colback=maincolor,
  sharp corners,
  boxrule=0pt,
  height=2.8cm,
  valign=center,
  before skip=-2.2cm,
  after skip=1.5cm,
  grow to left by=2.5cm,
  grow to right by=2.5cm,
]
  \begin{center}
    {\color{white}\sffamily\bfseries\Huge 斯特林公式的严谨推导}\\[3pt]
    {\color{white!85}\sffamily\large 基于伽马函数与拉普拉斯方法}
  \end{center}
\end{tcolorbox}

\vspace{0.5cm}

\begin{flushright}
  {\color{gray}\sffamily \today}
\end{flushright}

% ==================== 正文 ====================

\section{问题陈述}

\begin{problembox}
斯特林公式（Stirling's approximation）描述了当 $n \to \infty$ 时阶乘 $n!$ 的渐近行为。请利用\highlight{伽马函数}的积分表示，结合\highlight{拉普拉斯方法}，给出斯特林公式首项的完整且严谨的推导过程。
\[
  n! \sim \sqrt{2\pi n} \left( \frac{n}{e} \right)^n, \quad n \to \infty
\]
\end{problembox}

\section{推导：拉普拉斯方法}

\begin{methodbox}{核心思路}
通过变量代换 $x = nt$ 将移动峰值锚定在 $t=1$，分离大参数 $n$；在峰值处做二阶泰勒展开得到高斯近似；利用被积函数的快速衰减性将积分区间延拓至全实轴，最终由标准高斯积分引入 $\sqrt{2\pi}$ 常数。
\end{methodbox}

\subsection*{第一步：积分表示与变量代换}

由伽马函数定义出发：
\[
  n! = \Gamma(n+1) = \int_0^\infty x^n e^{-x} dx
\]
令 $x = nt$，则 $dx = n\,dt$，代入得：
\[
  n! = \int_0^\infty (nt)^n e^{-nt} \cdot n\,dt = n^{n+1} \int_0^\infty e^{n(\ln t - t)} dt
\]
此时积分具有拉普拉斯方法的标准形式 $\int e^{n\phi(t)} dt$，其中相函数为：
\[
  \phi(t) = \ln t - t
\]

\subsection*{第二步：相函数的极值分析}

对 $\phi(t)$ 求导寻找峰值：
\[
  \phi'(t) = \frac{1}{t} - 1, \quad \phi''(t) = -\frac{1}{t^2}
\]
令 $\phi'(t) = 0$ 得驻点 $t_0 = 1$。验证：
\begin{itemize}
  \item $\phi''(1) = -1 < 0$，确认 $t=1$ 为极大值点
  \item $\phi(1) = \ln 1 - 1 = -1$
\end{itemize}
变量代换的精妙之处在于：原本随 $n$ 移动的峰值 $x=n$ 被\highlight{锚定}在了固定的 $t=1$ 处。

\subsection*{第三步：局部泰勒展开与高阶项估计}

令 $t = 1 + u$，在 $u=0$ 附近展开 $\phi(1+u)$：
\[
  \phi(1+u) = -1 - \frac{u^2}{2} + \frac{u^3}{3} + O(u^4)
\]
代入积分并提取 $e^{-n}$：
\[
  n! = n^{n+1} e^{-n} \int_{-1}^\infty \exp\!\left(-\frac{n u^2}{2}\right) \cdot \exp\!\left(\frac{n u^3}{3} + O(n u^4)\right) du
\]

\begin{tcolorbox}[
  enhanced, colframe=accent!80, colback=accent!10,
  sharp corners, boxrule=1pt, left=4mm, right=4mm, top=3mm, bottom=3mm,
  before skip=0.8em, after skip=0.8em,
]
\sffamily\bfseries 为什么高阶项可以忽略？

\rmfamily\normalsize
有效积分区域宽度为 $u \sim O(n^{-1/2})$。主项 $nu^2 \sim O(1)$ 保留；三阶项 $nu^3 \sim O(n^{-1/2}) \to 0$；四阶项 $nu^4 \sim O(n^{-1}) \to 0$。因此 $\exp(nu^3/3 + \cdots) \to 1$，高阶项仅影响渐近级数的修正项，不影响首项近似。
\end{tcolorbox}

\subsection*{第四步：区间延拓与高斯积分}

简化后积分为：
\[
  n! \sim n^{n+1} e^{-n} \int_{-1}^\infty e^{-\frac{n u^2}{2}} du
\]
将下限从 $-1$ 替换为 $-\infty$。误差为 $\int_{-\infty}^{-1} e^{-nu^2/2} du = O(e^{-n/2})$，是\highlight{指数级小量}，相对于主积分的代数衰减 $O(n^{-1/2})$ 可完全忽略。

计算标准高斯积分，令 $v = u\sqrt{n}$：
\[
  \int_{-\infty}^\infty e^{-\frac{n u^2}{2}} du = \frac{1}{\sqrt{n}} \int_{-\infty}^\infty e^{-\frac{v^2}{2}} dv = \sqrt{\frac{2\pi}{n}}
\]
正是这个高斯积分引入了斯特林公式中关键的 $\sqrt{2\pi}$ 常数。

\subsection*{第五步：合并结果}

将所有因子组合：
\[
  n! \sim n^{n+1} e^{-n} \cdot \sqrt{\frac{2\pi}{n}} = \sqrt{2\pi n} \left( \frac{n}{e} \right)^n
\]

\newpage

\section{总结}

\begin{summarybox}

\subsection*{推导关键步骤回顾}

\begin{tabular}{@{}>{\sffamily\bfseries}p{3cm} p{5cm} p{5.5cm}@{}}
\toprule
\color{maincolor} \textbf{步骤} & \color{maincolor} \textbf{核心操作} & \color{maincolor} \textbf{数学意义} \\
\midrule
变量代换 & $x = nt$ & 将移动峰值固定为 $t=1$，分离大参数 $n$ \\
泰勒展开 & 保留至二阶 & 局部抛物线逼近，对应高斯近似 \\
区间延拓 & $[-1,\infty) \to (-\infty,\infty)$ & 利用被积函数快速衰减，误差指数级小 \\
高斯积分 & $\int e^{-v^2/2}dv = \sqrt{2\pi}$ & $\sqrt{2\pi}$ 常数的唯一来源 \\
\bottomrule
\end{tabular}

\vspace{0.5em}

\subsection*{进阶：斯特林渐近级数}

若不丢弃高阶项，系统展开可得完整渐近级数：
\[
  n! \sim \sqrt{2\pi n} \left(\frac{n}{e}\right)^n \left( 1 + \frac{1}{12n} + \frac{1}{288n^2} - \frac{139}{51840n^3} + O\!\left(\frac{1}{n^4}\right) \right)
\]
注意这是\highlight{渐近级数}而非收敛级数：对固定 $n$ 取过多项反而降低精度，但对固定项数当 $n \to \infty$ 时每项都提供更精确修正。

\subsection*{最终结论}

\begin{center}
\begin{tcolorbox}[
  enhanced, colframe=accent!80, colback=accent!10,
  sharp corners, boxrule=1.5pt, width=8cm, center,
]
  \centering
  {\color{maincolor}\sffamily\bfseries\LARGE 斯特林公式}\\[6pt]
  {\color{maincolor}\sffamily\bfseries\Large $\displaystyle n! \sim \sqrt{2\pi n} \left( \frac{n}{e} \right)^n$}
\end{tcolorbox}
\end{center}

\end{summarybox}

\end{document}